\documentclass[11pt,letterpaper]{article}
\usepackage[T1]{fontenc}
\usepackage[utf8]{inputenc}
\usepackage{lmodern}
\usepackage[margin=1in,headheight=15pt]{geometry}
\usepackage{amsmath,amssymb,amsthm,mathtools}
\usepackage{microtype,booktabs,array,longtable}
\newcolumntype{P}[1]{>{\raggedright\arraybackslash}p{#1}}
\usepackage[dvipsnames]{xcolor}
\usepackage{enumitem}
\usepackage{fancyhdr}
\usepackage{hyperref}
\usepackage{bookmark}
\hypersetup{colorlinks=true,linkcolor=MidnightBlue,citecolor=MidnightBlue,urlcolor=MidnightBlue,
 pdftitle={Surreal Scalars and Black-Hole Singularities},
 pdfauthor={Research assessment prepared for Vladimir Reshetnikov},
 pdfsubject={Surreal and surcomplex numbers in theoretical physics}}
\urlstyle{same}
\setlength{\emergencystretch}{2em}
\setlength{\parskip}{0.35em}
\setlength{\parindent}{1.1em}
\setlist{itemsep=3pt,topsep=5pt}
\pagestyle{fancy}
\fancyhf{}
\fancyhead[L]{\small Surreal scalars and black-hole singularities}
\fancyhead[R]{\small Research assessment}
\fancyfoot[C]{\thepage}
\renewcommand{\headrulewidth}{0.3pt}
\newcommand{\No}{\mathbf{No}}
\newcommand{\R}{\mathbb{R}}
\newcommand{\C}{\mathbb{C}}
\newcommand{\Q}{\mathbb{Q}}
\newcommand{\N}{\mathbb{N}}
\newcommand{\eps}{\varepsilon}
\newcommand{\st}{\operatorname{st}}
\newcommand{\ord}{\operatorname{ord}}
\newcommand{\supp}{\operatorname{supp}}
\newcommand{\Ric}{\operatorname{Ric}}
\newcommand{\diag}{\operatorname{diag}}
\newcommand{\dd}{\mathrm{d}}
\newcommand{\ii}{\mathrm{i}}
\newcommand{\Planck}{\ell_{\mathrm P}}
\newcommand{\Kretsch}{\mathcal K}
\newtheorem{proposition}{Proposition}[section]
\newtheorem{lemma}[proposition]{Lemma}
\newtheorem{corollary}[proposition]{Corollary}
\theoremstyle{definition}
\newtheorem{definition}[proposition]{Definition}
\newtheorem{example}[proposition]{Example}
\theoremstyle{remark}
\newtheorem{remark}[proposition]{Remark}
\newcommand{\status}[1]{\par\noindent\textbf{Status.} #1\par}

\begin{document}
\begin{titlepage}
\vspace*{0.5in}
\noindent{\small\sffamily MATHEMATICAL PHYSICS \quad / \quad FOUNDATIONS AND ASYMPTOTICS}
\vspace{0.45in}
\begin{flushleft}
{\LARGE\bfseries Surreal Scalars and\\[6pt] Black-Hole Singularities}\par
\vspace{0.24in}
{\Large What replacing $\R$ and $\C$ can accomplish,\\[3pt]
what it cannot accomplish by itself,\\[3pt]
and a rigorous research program}
\end{flushleft}
\vspace{0.35in}
\noindent Prepared for Vladimir Reshetnikov\par
\noindent 21 September 2026
\vspace{0.35in}
\begin{abstract}
The surreal field contains genuine infinite and infinitesimal numbers, and its
complexification is algebraically closed. These properties make it attractive
for describing the extreme hierarchies occurring near gravitational
singularities. They do not, however, turn a singular spacetime into a regular
one. We separate scalar extension, asymptotic representation, geometric
extension, and physical singularity resolution. Explicit Schwarzschild
calculations show how surreal evaluation retains invariant curvature blow-up,
finite-proper-time termination, and tidal deformation. A two-scale analysis of
radial versus nonradial infall illustrates a constructive use of ordered
asymptotic fields. We also prove a regular-reduction result for formal Einstein
equations, examine a nonsingular-core metric, and distinguish finite-dimensional
surcomplex quantum algebra from the analytic structures required for quantum
field theory. The assessment incorporates established surreal integration and
transseries results and the supplied repository's actual scope. The strongest
near-term role is a carefully specified asymptotic coefficient language over
ordinary spacetime, not an unqualified replacement of the continuum or a
stand-alone theory of quantum gravity.
\end{abstract}
\vfill
\noindent\small\textbf{Evidence and verification.} Literature and repository snapshot checked
as of 21 September 2026. This is a research assessment, not a claim of a new
empirically validated theory. The accompanying program checks 37 finite
symbolic identities. It does not verify the general analytic statements or the
repository's Lean build.
\end{titlepage}

\begingroup
\small
\setlength{\parskip}{0pt}
\tableofcontents
\endgroup
\newpage

\section{Executive assessment: a promising language, not a cure by arithmetic}
\label{sec:assessment}

There are two different questions hidden in the proposal to replace real and
complex numbers by surreal and surcomplex numbers.
The first is mathematical: can one formulate useful equations, asymptotic
expansions, geometries, or quantum models with these scalars? The answer is
\emph{yes}, provided the function classes, derivatives, sums, and interpretation
are specified. The second is physical: does the existence of infinite surreal
numbers itself remove black-hole singularities? The answer is \emph{no}. That
conclusion follows from explicit examples below, not from a general prejudice
against non-Archimedean models.

A singularity in general relativity is not normally a point of spacetime whose
coordinate has failed to fit inside $\R$. The basic singularity theorems concern
incomplete causal geodesics under geometric and causal hypotheses. Curvature
blow-up is a related but distinct issue. Replacing the scalar field does not,
by itself, supply an endpoint, an extension of the metric, a continuation law,
or an observational interpretation \cite{Penrose,Witten,Landsman}.

The genuinely constructive possibility is to represent processes \emph{approaching}
a singular regime by exact ordered scales. For example, if $\eps>0$ is
infinitesimal, then
\[
 \eps^{-6},\qquad \eps^{-2}\log(\eps^{-1}),\qquad
 \exp(\eps^{-1})
\]
are distinguishable, algebraically manipulable quantities. They carry more
information than one undifferentiated symbol $+\infty$. Established work on
surreal derivations, transseries, and integration gives substance to this
possibility; it is not merely a notational proposal
\cite{BerarducciMantova,ADH,CostinEhrlich}.

\begin{center}
\small
\begin{tabular}{@{}P{0.26\textwidth}P{0.63\textwidth}@{}}
\toprule
Proposed role & Assessment \\
\midrule
Exact representation of scales & Strong mathematical potential for selected algebraic,
transseries, and resurgent classes. A realization or summation theorem is still needed
when real functions are intended. \\
Perturbative tensor calculations & Feasible with real coordinate domains and controlled
formal or Hahn coefficient rings. Local algebraic identities survive. \\
Full replacement of spacetime scalars & A substantial foundational and analytic redesign,
not a substitution in existing formulas. \\
Automatic removal of black-hole singularities & Not established and refuted as a general
mechanism by invariant-pole examples. \\
New fundamental physical theory & Logically possible, but requires dynamics, measurement,
recovery of known physics, and discriminating predictions beyond scalar extension. \\
\bottomrule
\end{tabular}
\end{center}

The remainder develops this assessment constructively. Some propositions are
included to make the logical boundaries explicit. They are elementary
consequences of the stated assumptions or standard calculations presented here
in full; no claim of priority is made for them. Open questions are distinguished
from facts proved in the article.

\section{What is being enlarged?}
\label{sec:scalars}

\subsection{Surreal and surcomplex numbers}

The surreal numbers $\No$ form a real-closed ordered field in the appropriate
class-sized sense. They contain the real numbers, infinite elements such as
$\omega$, and positive infinitesimals such as $\omega^{-1}$. Their
complexification
\[
 \No[\ii]=\{a+\ii b:a,b\in\No\},\qquad \ii^2=-1,
\]
is algebraically closed. Conjugation and modulus are defined by
\[
 \overline{a+\ii b}=a-\ii b,
 \qquad |a+\ii b|=\sqrt{a^2+b^2}\in\No_{\geq 0}.
\]
The latter is a surreal-valued modulus, not an ordinary real-valued Banach-space
norm \cite{Conway,Gonshor}.

These algebraic facts preserve much finite-dimensional mathematics. One can
solve polynomial equations, invert nonsingular matrices, construct quadratic
forms, and manipulate tensor components. None of these facts alone supplies
a satisfactory theory of smooth functions, partial differential equations,
measures, or infinite-dimensional Hilbert spaces.

\begin{definition}
In an ordered field $F\supseteq\R$, an element $x$ is \emph{limited} if
$|x|<M$ for some real $M>0$. It is \emph{infinitesimal} if $|x|<r$ for every
real $r>0$. An element not limited is \emph{unlimited}.
\end{definition}

This vocabulary prevents an important ambiguity. An infinite surreal is an
individual field element, but it is not thereby a finite physical magnitude
relative to ordinary units. Saying that curvature has become ``a number'' is
not the same as saying that a curvature observable is limited.

Moreover, $\eps\neq0$ does not make division by zero legal. In every nontrivial
field, including $\No$, the equation $0y=1$ has no solution. An infinite surreal
has a nonzero infinitesimal inverse; it is not the reciprocal of zero. Nor do
surreals provide a smallest positive length: $0<\eps/2<\eps$. The Planck length,
when treated as an ordinary physical constant, is a small positive real length,
not a literal positive infinitesimal.

\subsection{Normal forms and set-sized workspaces}

A surreal normal form has the schematic shape
\[
 x=\sum_{\alpha<\lambda}r_\alpha\omega^{a_\alpha},
 \qquad r_\alpha\in\R\setminus\{0\},
 \qquad a_\alpha\ \text{strictly decreasing},
\]
with \emph{set-sized} support. This is not a sum over all ordinals. For most
calculations it is better to fix a set-sized ordered abelian group $\Gamma$ and
work in a Hahn field
\[
 F_\Gamma=\R((t^\Gamma)),\qquad
 K_\Gamma=\C((t^\Gamma)).
\]
Its elements have well-ordered exponent support, and a positive exponent
represents a smaller scale. If $\Gamma$ is divisible, these are respectively
real closed and algebraically closed. An ordered additive embedding of
$\Gamma$ into surreal exponents identifies such series with an appropriate
surreal subfield using $t^\gamma\mapsto\omega^{-\gamma}$, with the embedding
understood in the exponent \cite{Conway,Gonshor,RepoDocs}.

Any concrete set of surreal coefficients has only a set of exponents in its
combined support. Taking the divisible group generated by those exponents
produces a workspace containing that data. This observation removes a
proper-class obstacle from individual calculations. It does \emph{not} prove
that an arbitrary chosen workspace is closed under every desired derivative,
exponential, integral, or infinite summation. Those closure properties must be
specified separately.

The full proper class is therefore not an obvious physical advantage. It is
an ambient language within which particular set-sized structures can be
selected. For computation and proof assistance, the smaller structure is
usually the meaningful object.

\subsection{What summation means}

A family of Hahn series is \emph{strongly summable} if the union of its supports
is well ordered and each exponent receives contributions from only finitely
many members. These are algebraic support conditions. They are different from
ordinary convergence of partial sums \cite{Gonshor,RepoAnalysis}.

For instance, with $\eps=\omega^{-1}$,
\[
 \sum_{n\geq0}\eps^n=(1-\eps)^{-1}
\]
is a valid Hahn identity. Its finite remainder identity is simply
\begin{equation}
 (1-\eps)^{-1}-\sum_{n=0}^{N}\eps^n
   =\frac{\eps^{N+1}}{1-\eps}.
 \label{eq:geom}
\end{equation}
This equation is useful without claiming convergence in the full surreal order
topology. Similarly, $\sum_{n\geq0}n!\eps^{n+1}$ is strongly summable: its
support is a subset of the positive integers and each exponent occurs once.
The corresponding real power series has zero radius of convergence. The
surreal sum does not automatically select a particular real function with
that asymptotic expansion; Section~\ref{sec:resurgence} makes the distinction
explicit.

\section{The analytic structures that do not transfer automatically}
\label{sec:analysis}

\subsection{The full order topology is unsuitable for ordinary histories}

Use order balls on $\No$, or modulus balls on $\No[\ii]$, as the default
neighborhoods. The following observation explains why a direct substitution
in standard real analysis is unsafe.

\begin{proposition}[Set-sized discreteness]
\label{prop:discrete}
Every set-sized subset of $\No[\ii]$ is closed and discrete in this topology.
Every convergent set-indexed net is eventually equal to its limit. Every
continuous map from an ordinary connected real interval into $\No[\ii]$ is
constant.
\end{proposition}
\begin{proof}
Let $A$ be a set of surcomplex numbers and let $x$ be any surcomplex point.
The nonzero distances $|x-a|$, for $a\in A\setminus\{x\}$, form a set of
positive surreals. The surreal cut construction gives a positive $\delta$
smaller than every member of that set. Thus the ball of radius $\delta$ at
$x$ meets $A$ only at $x$, if $x\in A$, and otherwise misses $A$.
This proves discreteness and closedness. Empty sets of distances cause no
difficulty: any positive radius may be chosen.

For a set-indexed net, its range together with a proposed limit is a set.
A neighborhood of the limit isolates it from every other point in this set;
convergence forces eventual equality. Finally, the image of a real interval
is a set, hence discrete, whereas a continuous image of a connected space is
connected. The image therefore contains at most one point.
\end{proof}

In particular, the ordinary inclusion of a real interval in the full surreal
field is not continuous for this finer topology. Even ordinary real sequences
converging in $\R$ need not converge after inclusion. The proposition does not
say that the full class is a discrete space: a singleton is not an open ball,
and the proof cannot choose one radius below a \emph{proper class} of positive
distances.

There are several ways around this obstacle, but each changes the model. One
can use a valuation topology on a set-sized field, a coefficientwise topology,
a standard-part topology, or a specialized notion of surreal analytic function.
For a rank-one Hahn field, powers $t^n$ can tend to zero in its own valuation
topology, even though this convergence is destroyed by inclusion in the full
surreal order topology. One can also retain ordinary real spacetime as the
parameter domain and attach formal coefficients to its fields. That last
option will be our constructive framework.

The supplied repository emphasizes precisely the distinction between Hahn
summation and topological convergence. It also distinguishes functions given
by local germs, coherent coefficient functions on a common ordinary domain,
and lifts of ordinary holomorphic functions. A theorem for one class cannot
be silently applied to another \cite{RepoDocs,RepoAnalysis}.

\subsection{Three different meanings of differentiation}

The following operations answer different questions.

\paragraph{Spacetime differentiation.}
For a field $f(x;\eps)=\sum_n f_n(x)\eps^n$ on an ordinary real coordinate
patch, one may define
\[
 \partial_\mu f=\sum_n (\partial_\mu f_n)\eps^n,
 \qquad \partial_\mu\eps=0.
\]
The small parameter is a scalar constant. This is the derivative required for
formal metric perturbations.

\paragraph{A derivation of the scalar field.}
Berarducci and Mantova construct a derivation on $\No$ compatible with its
exponential and asymptotic structure. In its usual normalization $D\omega=1$,
so $D(\omega^{-1})=-\omega^{-2}$, whereas real constants are annihilated.
Subsequent work relates this differential field to transseries and Hardy fields
\cite{BerarducciMantova,ADH}. This is valuable asymptotic differential algebra.
It is not automatically a physical time derivative. In particular, applying
$D$ to a real number used as a time coordinate does not recover differentiation
of a real time-dependent function.

\paragraph{Differentiation of surreal functions.}
A function on a surreal domain can have its own derivative, subject to a
specified function theory. Costin and Ehrlich develop extensions and integration
for substantial classes associated with resurgent analysis, including suitable
ordinary differential and difference equations \cite{CostinEhrlich}. This is
positive evidence that meaningful surreal calculus exists. It does not assert
a general extension theorem for arbitrary smooth spacetime fields or the full
nonlinear Einstein initial-value problem.

These distinctions matter physically. Choosing a scalar derivation by analogy
with time evolution can inadvertently make a supposedly constant regulator
time-dependent. Conversely, declaring all infinitesimals constant is compatible
with coefficientwise perturbation theory but does not capture every
asymptotic differential-field construction. A model must state which operation
is used and prove its chain rules in the chosen category.

\subsection{Transfer: substantial, but language-dependent}

Surreals do not lack all transfer principles. Real-closed fields satisfy the
same first-order theory in the language of ordered fields. There are also
strong elementary-extension results for surreal exponentiation
\cite{vanDenDriesEhrlich}. These results justify many algebraic and definable
arguments.

They do not transfer arbitrary theorems about real functions, compact sets,
Lebesgue integration, distributions, or Hilbert spaces. Those objects require
additional sorts, predicates, or constructions. Robinson-style nonstandard
analysis packages a transferred mathematical universe with distinguished
internal objects; an ordered-field embedding into the surreals does not by
itself retain that entire package. It would be equally misleading to say
``everything transfers'' or ``nothing transfers.'' The question is which
language, structure, and hypotheses have actually been transported.

\section{What a black-hole singularity is}
\label{sec:singularity}

\subsection{Four notions that must remain separate}

A coordinate formula can diverge even when the underlying geometry is regular.
A curvature tensor can diverge in a freely falling frame even if some scalar
invariants remain bounded. A causal geodesic can be incomplete even without
an established curvature blow-up theorem. A spacetime can admit a continuous
metric extension but no extension with the differentiability needed for the
classical field equations. These are different defects, requiring different
notions of resolution \cite{Witten,Landsman,DafermosLuk}.

For present purposes, a genuine resolution must specify what replaces the
incomplete region and what governs propagation there. Depending on the theory,
this might mean a regular Lorentzian extension, a weaker geometry with a
well-posed evolution law, or a quantum description in which operational
observables and evolution remain meaningful. It need not mean that every
classical variable stays finite, but then the replacement observables and
rules must be stated explicitly.

The phrase ``usual real analysis fails'' therefore needs refinement. Real
analysis often succeeds at proving that a particular classical solution cannot
be extended with a specified regularity. That is a diagnosis of the solution
or model, not a contradiction in the real numbers. A different number system
can help express the diagnosis or organize a replacement model. It does not
invalidate the diagnosis without changing the relevant assumptions.

\subsection{The singularity theorem is not a division-by-zero theorem}

One standard form of Penrose's theorem combines a trapped surface, the null
convergence condition, and suitable global causal assumptions including a
noncompact Cauchy surface to infer future null geodesic incompleteness
\cite{Penrose,Witten}. Its proof combines local focusing with global geometry.
It does not merely encounter a fraction with a vanishing denominator, and it
does not assert that every incomplete geodesic has divergent Kretschmann
scalar.

The local null Raychaudhuri equation, in four dimensions for a twist-free
congruence, gives under the null convergence condition
\begin{equation}
 \frac{\dd\theta}{\dd\lambda}
 =-\frac12\theta^2-\sigma_{ab}\sigma^{ab}
   -R_{ab}k^ak^b
 \leq-\frac12\theta^2.
 \label{eq:ray}
\end{equation}
An initially negative expansion is driven to focusing within affine parameter
at most $2/|\theta_0|$, as long as the regular evolution in question exists.
The algebraic inequality can be represented in an ordered extension field.
The global assumptions concerning causality and completeness cannot simply be
replaced by field arithmetic. Conversely, showing that the classical proof
cannot be copied into an unfamiliar topology does not exhibit a nonsingular
physical spacetime.

\subsection{Affine and proper time are not optional conventions}

Suppose a timelike geodesic reaches the end of its domain at finite proper time
$\tau_*$. Introducing $u=-\log(\tau_*-\tau)$ moves that endpoint to $u=+\infty$.
This does not make the geodesic complete: $u$ is not its proper-time parameter,
and in general is not affine. Incompleteness is not cured by relabeling the
clock. A surreal value for $u$ may encode an extraordinarily late stage in
this reparameterization, but the missing proper-time continuation remains a
separate problem.

\section{Schwarzschild: what infinitesimal radius actually represents}
\label{sec:schwarzschild}

We use signature $(-,+,+,+)$ and geometrized units $G=c=1$ unless otherwise
stated. Thus mass $m=GM/c^2$, radial distance, and proper time expressed as a
length have the same dimension. Physical seconds are recovered by dividing
proper-time lengths by $c$.

\subsection{The horizon is already regular over the reals}

The Schwarzschild metric in its familiar static chart is
\begin{equation}
 \dd s^2=-f(r)\dd t^2+f(r)^{-1}\dd r^2+r^2\dd\Omega^2,
 \qquad f(r)=1-\frac{2m}{r}.
 \label{eq:schwarzmetric}
\end{equation}
The coefficient $f^{-1}$ diverges at $r=2m$. Introduce the advanced coordinate
$v=t+r_*$, where $\dd r_*/\dd r=f^{-1}$. Direct substitution gives
\begin{equation}
 \dd s^2=-f(r)\dd v^2+2\dd v\dd r+r^2\dd\Omega^2.
 \label{eq:ef}
\end{equation}
Its determinant is $-r^4\sin^2\theta$, nonzero at $r=2m$ away from the ordinary
angular-coordinate poles. Thus the apparent horizon divergence did not require
new scalars. A different real coordinate chart removes it.

For any metric of the form \eqref{eq:schwarzmetric}, a direct curvature
calculation yields
\begin{equation}
 \Kretsch:=R_{abcd}R^{abcd}
   =[f''(r)]^2+4\left[\frac{f'(r)}r\right]^2
       +4\left[\frac{1-f(r)}{r^2}\right]^2.
 \label{eq:generalK}
\end{equation}
Consequently Schwarzschild has
\begin{equation}
 \Kretsch=\frac{48m^2}{r^6},
 \qquad \Kretsch\big|_{r=2m}=\frac{3}{4m^4}.
 \label{eq:schwarzK}
\end{equation}
The companion program derives \eqref{eq:generalK} by contracting the full
Riemann tensor, rather than merely substituting into an assumed Schwarzschild
answer. It also verifies $R_{ab}=0$ and the determinant in
\eqref{eq:ef}. The physical distinction between a regular horizon and the
central singularity is standard \cite{Witten}.

\subsection{The center is not fixed by evaluating at a small nonzero number}

Let $L>0$ be an ordinary reference length, set $m=\mu L$ with $\mu>0$ real,
and evaluate the exterior/interior formula at a positive infinitesimal radius
$r=L\eps$. Then
\begin{equation}
 L^4\Kretsch=48\mu^2\eps^{-6}.
 \label{eq:infK}
\end{equation}
This is a well-defined, unlimited surreal number. It identifies the exact
order of the curvature divergence. It does not make that divergence disappear.
The statement is dimensionally meaningful because the compared quantity is
$L^4\Kretsch$, not a bare dimensionful number declared ``infinite'' without
units.

The radius here is not an arbitrary coordinate label: it is the areal radius,
defined by the area $4\pi r^2$ of symmetry spheres. A change of coordinates
cannot change the scalar value in \eqref{eq:schwarzK} at a fixed event.
One may alter the notion of event or replace the spacetime model, but that is
more than an algebraic relabeling.

\begin{proposition}[A leading pole stays unlimited]
\label{prop:pole}
Let $F\supseteq\R$ be an ordered field, let $\eps>0$ be infinitesimal, and let
$p>0$ be such that $\eps^p$ is defined and infinitesimal in the chosen field.
If $a\in\R\setminus\{0\}$ and $u$ is infinitesimal, then
\[
 I=a\eps^{-p}(1+u)
\]
is unlimited. An ordered extension of $F$ that fixes the real numbers does not
make this same element limited.
\end{proposition}
\begin{proof}
Since $|u|<1/2$, $|I|>|a|\eps^{-p}/2$. The latter exceeds every positive real
bound by infinitesimality of $\eps^p$. Every inequality with a real bound is
preserved by an ordered extension.
\end{proof}

Thus changing from a smaller ordered asymptotic field to all of $\No$ cannot
turn the same curvature observable into a limited one. Subtracting terms,
rescaling the measured observable by an infinitesimal, or changing the metric
can change the conclusion, but each introduces additional structure or a new
physical claim.

There is a second, purely algebraic obstruction to filling in the original
formula at the exact center. On the Schwarzschild region,
\begin{equation}
 r^6\Kretsch=48m^2.
 \label{eq:polyobstruction}
\end{equation}
No field-valued assignment with $r=0$ and $m\ne0$ can satisfy this identity.
A regular extension preserving this relation at its endpoint is therefore
impossible in any field, not just $\R$. This is a conditional obstruction to
that kind of extension, not a theorem excluding every proposed quantum or
nonclassical replacement of the endpoint. Setting $r=\eps$ merely avoids
$r=0$; it does not provide the missing assignment.

\subsection{Exact radial proper-time asymptotics}

For a radial timelike geodesic, the conserved specific energy is
$E=f\,\dd t/\dd\tau$. Normalization of the four-velocity gives
\[
 \left(\frac{\dd r}{\dd\tau}\right)^2=E^2-f(r).
\]
Choose infall with $E=1$, and let
\[
 s=\tau_*-\tau>0
\]
be the remaining proper-time length before the classical endpoint. Then
\[
 \frac{\dd r}{\dd\tau}=-\sqrt{\frac{2m}{r}},
 \qquad
 r^{3/2}=\frac32\sqrt{2m}\,s.
\]
Therefore
\begin{equation}
 r=A s^{2/3},\qquad A=\left(\frac{9m}{2}\right)^{1/3},
 \qquad
 \frac{m}{r^3}=\frac{2}{9s^2},\qquad
 \Kretsch=\frac{64}{27s^4}.
 \label{eq:radial}
\end{equation}
These formulas are exact for this geodesic family. They are especially suited
to a Puiseux or Hahn representation. Taking $s=L\eps$ produces exact
infinitesimal distances and unlimited curvature while retaining the relation
to the infaller's clock.

The endpoint is still at finite proper time. Replacing a real $s>0$ by a
positive infinitesimal specifies an extraordinarily close approach in the
enlarged model. It does not define the solution at $s=0$ or decide what follows
it. Allowing negative $s$ in the algebraic relation
$r^3=(9m/2)s^2$ produces a formal other branch, but that relation alone omits
the differentiability, curvature, orientation, and matching information needed
for a physical extension.

\subsection{Tidal deformation is not hidden by scalar extension}

In a freely falling orthonormal frame along radial Schwarzschild infall, the
radial and transverse Jacobi equations take the form
\begin{equation}
 \frac{\dd^2J_r}{\dd\tau^2}=\frac{2m}{r^3}J_r,
 \qquad
 \frac{\dd^2J_\perp}{\dd\tau^2}=-\frac{m}{r^3}J_\perp.
 \label{eq:jacobi}
\end{equation}
Using $\dd^2/\dd\tau^2=\dd^2/\dd s^2$ and \eqref{eq:radial}, the indicial
equations are
\[
 p(p-1)=\frac49 \quad\text{and}\quad p(p-1)=-\frac29.
\]
Thus
\begin{align}
 J_r&=C_1s^{4/3}+C_2s^{-1/3},\label{eq:Jr}\\
 J_\perp&=D_1s^{2/3}+D_2s^{1/3}.\label{eq:Jt}
\end{align}
For generic nonzero leading coefficients, radial separation grows like
$s^{-1/3}$ while two independent transverse separations shrink like $s^{1/3}$.
The corresponding infinitesimal volume scales like $s^{1/3}$ and tends to zero.
Special initial conditions can remove leading terms and change these rates;
the displayed generic behavior is not a claim about every Jacobi field.

A surreal evaluation records both the unlimited stretching and the
infinitesimal collapse. Calling each value an available field element does not
supply a surviving detector, restore a nonzero volume, or determine propagation
through the endpoint. This example is more informative than the statement
``infinite curvature is now allowed'': it connects the singular regime with
an observer's clock and neighboring freely falling trajectories.

\section{A constructive multiscale application: small angular momentum}
\label{sec:angular}

The preceding $s^{-4}$ curvature law should not be universalized to every
Schwarzschild geodesic. The dependence on angular momentum gives a useful
example of what an ordered asymptotic field can genuinely organize.

\subsection{Nonradial infall changes the proper-time exponent}

Restrict by spherical symmetry to the equatorial plane and let
$j=r^2\dd\phi/\dd\tau$ be the conserved specific angular momentum. For finite
$E$ and $j\ne0$,
\begin{equation}
 \dot r^{\,2}=E^2-
    \left(1-\frac{2m}{r}\right)\left(1+\frac{j^2}{r^2}\right)
 =\frac{2mj^2}{r^3}+\frac{2m}{r}-\frac{j^2}{r^2}+E^2-1.
 \label{eq:radialgeneral}
\end{equation}
The first term dominates as $r\downarrow0$ with $j$ fixed and nonzero.
Integrating the leading balance yields
\begin{equation}
 r\sim\left[\frac52\sqrt{2m}\,|j|\,s\right]^{2/5},
 \quad
 \Kretsch\sim
 48m^2\left[\frac52\sqrt{2m}\,|j|\right]^{-12/5}s^{-12/5}.
 \label{eq:nonradial}
\end{equation}
The divergence remains, but its power in remaining proper time changes. This
is not a change in the curvature as a function of areal radius. It is a change
in how the worldline approaches the same invariantly characterized region.

\subsection{A uniform crossover problem}

Let $J=|j|>0$, set $r=JR$, and define the dimensionless remaining clock
\[
 S=\frac{\sqrt{2m}}{J^{3/2}}s.
\]
For the ingoing branch that reaches $r=0$, integration of the exact first
integral gives
\begin{equation}
 S=\int_0^R
 \frac{u^{3/2}\,\dd u}
 {\sqrt{1+u^2+\dfrac{J}{2m}\bigl[(E^2-1)u^3-u\bigr]}}.
 \label{eq:crossoverexact}
\end{equation}
The radicand is positive in the sufficiently small interior region under
consideration. For $J/m\to0$ with $E$ fixed and $R$ in any bounded positive
interval, its reciprocal square root converges uniformly to that of $1+u^2$.
Consequently
\begin{equation}
 S\longrightarrow\Phi(R):=
   \int_0^R\frac{u^{3/2}}{\sqrt{1+u^2}}\,\dd u.
 \label{eq:Phi}
\end{equation}
The elementary endpoint expansions are
\begin{equation}
 \Phi(R)\sim\frac25R^{5/2}\quad(R\downarrow0),
 \qquad
 \Phi(R)\sim\frac23R^{3/2}\quad(R\to\infty).
 \label{eq:Philimits}
\end{equation}
The first yields the angular-momentum-dominated law; the second recovers the
radial law in the overlap region $J\ll r\ll m$. The large-$R$ asymptotic of the
limiting function is a matched-asymptotic statement: it is not a claim that
convergence in \eqref{eq:Phi} is uniform on an unbounded interval.

This construction uses ordinary real analysis to supply an actual crossover
function. Hahn valuations then give a concise way to select its asymptotic
regions. With $m=\mu L$, $J=L\eps^\beta$, and $s=L\eps^\alpha$, where
$\alpha,\beta>0$ are real, the consistency conditions are
\begin{align}
 \alpha<\frac32\beta &: \quad
        r/L\ \text{has leading exponent }\frac{2\alpha}{3},
        \qquad r\gg J,\label{eq:radialsector}\\
 \alpha>\frac32\beta &: \quad
        r/L\ \text{has leading exponent }\frac{2(\alpha+\beta)}{5},
        \qquad r\ll J.\label{eq:angularsector}
\end{align}
At equality neither monomial should be discarded; $\Phi$ describes the
transition. Real coefficients of order one can be restored without changing
these exponent tests.

This is an example of useful surreal-compatible physics mathematics: retain
several small scales, determine which terms dominate, and avoid taking
incompatible limits. It offers no new gravitational law and is equally
expressible in suitable transseries or generalized-series fields. Its value
must be judged by clarity, support control, and computational reliability,
not by whether the ambient field contains a proper class of still larger
numbers.

\section{The classical theory loses control before a literal infinitesimal scale}
\label{sec:EFT}

A mathematical extension of a classical formula and a controlled physical
prediction are different things. In a low-energy effective description of
gravity, higher-curvature and quantum corrections are organized relative to
appropriate scales. The coefficients, field content, and state are physical
inputs, not consequences of choosing $\No$ \cite{Donoghue}.

Restore
\[
 \Planck=\sqrt{\frac{\hbar G}{c^3}},\qquad m=\frac{GM}{c^2}.
\]
A dimensional indicator for large curvature corrections is
$\Planck^4\Kretsch$. For Schwarzschild, setting this indicator to order one
gives
\begin{equation}
 r_{\mathrm Q}=48^{1/6}(m\Planck^2)^{1/3}.
 \label{eq:rQ}
\end{equation}
This is an estimate of a curvature scale, not a derived minimum radius or a
universal theorem identifying where a particular quantum-gravity theory takes
over. Additional scales or large state-dependent effects can invalidate a
simple curvature criterion sooner.

For $m\gg\Planck$, the ratio $r_{\mathrm Q}/(2m)$ is small, while
$r_{\mathrm Q}/\Planck$ is large. Thus a macroscopic areal radius inside a large
black hole can coexist with Planck-scale curvature. ``Small radius'' and
``high curvature'' cannot be identified without specifying the mass.

Now fix ordinary positive real $L$ and $\Planck/L$. A literal positive
infinitesimal $\eps$ satisfies $\eps<\Planck/L$, and indeed lies below every
other positive real cutoff ratio. Evaluating the unchanged classical solution
at $r=L\eps$ explores a regime beyond any such fixed cutoff. It does not
justify ignoring the terms that made the classical approximation unreliable.

An ordered asymptotic field can instead encode a \emph{joint family}, for
example $\Planck/L=\eps^b$ and $r/L=\eps^a$. That can be useful for comparing
classical and quantum corrections symbolically. But it is a formal scaling
model, not a deduction that the physical Planck length is itself infinitesimal.
In a higher-curvature effective action, one must still specify the independent
operators and their coefficients; in four dimensions some quadratic curvature
terms are related by identities or field redefinitions. Merely writing
additional powers of curvature does not determine a UV completion.

\section{Regularization, extension, and new dynamics}
\label{sec:regularization}

\subsection{A pole in an ordinary differential equation}

Consider
\[
 \frac{\dd y}{\dd t}=y^2,\qquad y(t)=\frac{1}{t_*-t}.
\]
At $t=t_*-\eps$, the solution is $y=\eps^{-1}$. At $t=t_*$, the identity
$(t_*-t)y=1$ has no field-valued solution. This is the simplest illustration of
the difference between an infinite nearby value and a value at the pole.

There is nevertheless a legitimate geometric maneuver: introduce the
projective coordinate $z=1/y$. Then $z=t_*-t$ is smooth through zero, and the
vector field becomes $\dot z=-1$. A compactified state space may therefore
continue the \emph{projective} trajectory. Whether this solves a physical
problem depends on whether the physical observables and evolution laws are
regular in that state space. If $y$ itself is a measured stress, replacing it
by $z$ has not made the measured stress bounded.

This example guards against two opposite errors. Not every divergence is a
fundamental obstruction, because some are removable by changing the geometric
description. But not every compactification is a physical resolution, because
the regularity of the relevant observables must be checked. Schwarzschild's
scalar curvature pole is not removed merely by complexifying a coordinate or
adjoining an infinite scalar.

\subsection{Complex continuation is not causal continuation}

Surcomplex values can be useful in algebraic or analytic continuation, just
as ordinary complex values already are. For instance, fractional powers in
\eqref{eq:radial} admit complex branches. A path in a complexified coordinate
domain may avoid a pole or branch point.

To turn such a construction into a real physical continuation, one must
recover a real Lorentzian slice, specify the reality conditions, choose the
branch and matching data, and show that the resulting evolution satisfies the
intended equations. Neither the ordinary complex plane nor its surcomplex
enlargement supplies these choices. Analytic continuation in a parameter is
not the same operation as extending an infalling observer's causal worldline.

\subsection{What a proposed singularity resolution must deliver}

A useful test is to require an explicit answer to four questions. What are the
states or events replacing the classical endpoint? What equations govern
propagation there? Which observables remain meaningful, and in what sense?
Does the replacement determine a continuation uniquely from admissible data,
or does it introduce new boundary conditions?

The answer may involve nonclassical geometry or quantum states rather than a
smooth metric. That is allowed. What is not enough is to retain every classical
formula, replace the word ``divergent'' by ``surreal-valued,'' and regard the
missing choices as settled.

\section{A rigorous coefficientwise formulation of gravitational perturbations}
\label{sec:formalGR}

The most conservative useful architecture keeps spacetime real and enlarges
its coefficient algebra. This avoids pretending that ordinary real histories
are continuous in the full surreal order topology.

\subsection{The differential algebra}

Let $U$ be a real coordinate domain and set
\[
 \mathcal A(U)=C^\infty(U)[[\eps]].
\]
Coefficients are ordinary smooth functions on \emph{one common domain}. Products
are formal Cauchy products, with only finitely many contributions at each
integer order. Define differentiation coefficientwise and let
\[
 \pi_0:\mathcal A(U)\longrightarrow C^\infty(U),
 \qquad \sum_{n\geq0}f_n\eps^n\longmapsto f_0.
\]
Then $\pi_0$ is a ring homomorphism and
$\pi_0\partial_\mu=\partial_\mu\pi_0$.

This is initially formal differential algebra. At an ordinary point $x$, the
series can be evaluated in an appropriate Hahn field, but that fact alone does
not make the map $x\mapsto f(x;\eps)$ smooth in the full surreal order topology.
The specified derivative, not an unstated topological identification, carries
the geometric construction.

More elaborate rings $C^\infty(U)((t^\Gamma))$ can be useful, but require common
domains, controlled supports, and summability under the operations employed.
A collection of coefficient germs with shrinking domains is not the same
object. The repository's explicit separation of these coefficient categories
is directly relevant here \cite{RepoDocs}.

\subsection{Formal metric inversion and curvature}

Let
\begin{equation}
 g_{\mu\nu}=g^{(0)}_{\mu\nu}
       +\eps g^{(1)}_{\mu\nu}+\eps^2g^{(2)}_{\mu\nu}+\cdots,
 \label{eq:formalg}
\end{equation}
where $g^{(0)}$ is a smooth nondegenerate Lorentzian metric on $U$. Since
$\det g^{(0)}$ never vanishes, $g$ has a formal inverse. In matrix notation,
write $g=g^{(0)}(I+H)$ with $H$ of positive order. Then
\begin{equation}
 g^{-1}=\left(\sum_{n\geq0}(-H)^n\right)(g^{(0)})^{-1}.
 \label{eq:ginv}
\end{equation}
Every coefficient is a finite sum, so this formula needs no analytic
convergence assumption.

Define the Levi-Civita coefficients by their usual finite expression,
\[
 \Gamma^\rho_{\mu\nu}
 =\frac12g^{\rho\sigma}
 (\partial_\mu g_{\sigma\nu}+\partial_\nu g_{\sigma\mu}
                  -\partial_\sigma g_{\mu\nu}),
\]
and then form the curvature and Einstein tensor. These constructions make
sense in the differential algebra. The usual tensor identities and contracted
Bianchi identity hold coefficientwise because their proofs use the same
product rule, commuting coordinate derivatives, inverse identities, and finite
contractions. They do not assert global existence, stability, or causal
properties of a spacetime over a new scalar field.

\begin{proposition}[Regular reduction of the Einstein equations]
\label{prop:reduction}
Let $g$ be as in \eqref{eq:formalg} and let $T$ be a formal symmetric tensor.
For ordinary real constants $\Lambda$ and $\kappa$, suppose
\begin{equation}
 G[g]+\Lambda g=\kappa T
 \label{eq:formalEinstein}
\end{equation}
holds coefficientwise on $U$. Then
\begin{equation}
 G[g^{(0)}]+\Lambda g^{(0)}=\kappa T^{(0)}.
 \label{eq:shadowEinstein}
\end{equation}
In particular, $\pi_0(G[g])=G[\pi_0(g)]$.
\end{proposition}
\begin{proof}
Reduction preserves products and sums, commutes with coordinate derivatives,
and sends the inverse in \eqref{eq:ginv} to $(g^{(0)})^{-1}$. Apply these facts
successively to the connection, Riemann tensor, Ricci tensor, scalar curvature,
and Einstein tensor. Reducing \eqref{eq:formalEinstein} gives
\eqref{eq:shadowEinstein}.
\end{proof}

This gives a precise consistency theorem: ordinary general relativity is
recovered at leading order in a regular formal deformation. Higher orders
encode perturbations; at first order one obtains the linearized Einstein
operator together with the corresponding matter perturbation.

The same theorem limits what this regular construction can claim. If the
leading metric has a classical singular boundary, the theorem does not extend
it through that boundary. Its hypotheses require a common smooth domain and
an invertible leading metric. A proposed resolution may involve an inner
region where those hypotheses fail, resummed terms become dominant, or a new
leading geometry appears. That is a singular perturbation problem, not a
contradiction of the reduction theorem.

\subsection{Gauge transformations and support control}

Ordinary coordinate transformations of $U$ act on the coefficient tensors in
the usual way. Formal diffeomorphisms near the identity can also be defined
order by order using Lie derivatives. Neither operation changes the invariant
content merely because the coefficients lie in a larger algebra.

A practical implementation should treat support conditions as part of every
operation's input contract. Inversion requires an invertible leading term;
composition requires the relevant positive-order or domain hypotheses;
integration requires both a coefficientwise integral and an admissible combined
support. It is unsafe to derive an identity in a formal ring, interpret it as
an ordinary convergent sum, and then use ordinary interchange-of-limit theorems
without proving a bridge.

\section{Measurement, standard part, and noncanonical subtraction}
\label{sec:measurement}

\subsection{The standard-part map has a precise domain}

For an ordered field $F\supseteq\R$, write
\[
 \mathcal O_F=\{x\in F:x\text{ is limited}\},\qquad
 \mathfrak m_F=\{x\in F:x\text{ is infinitesimal}\}.
\]
Real completeness implies that every limited $x$ is infinitely close to a
unique real number. Indeed, take the supremum in $\R$ of the real numbers
below $x$; any noninfinitesimal discrepancy would contradict the defining
upper-bound property. This defines
\[
 \st:\mathcal O_F\longrightarrow\R,
 \qquad \ker\st=\mathfrak m_F.
\]
The map is a ring homomorphism, and
$\mathcal O_F/\mathfrak m_F\cong\R$. Its complex version applies separately to
real and imaginary parts. The repository implements a corresponding
coefficient-zero reduction for its nonnegative-order Hahn subring
\cite{RepoRoot}.

A model can postulate that limited quantities are read by ordinary apparatus
through standard part. That is an intelligible correspondence rule, but it is
an additional modeling assumption. It is not defined on an unlimited
curvature such as \eqref{eq:infK}. Extending the codomain to include
$+\infty$ still does not provide an ordinary finite detector output.

The reduction theorem in Section~\ref{sec:formalGR} used a particularly strong
compatibility: reduction commuted with differentiation \emph{by definition} in
a regular formal algebra. This does not hold for arbitrary families followed
by a limiting or standard-part operation. For example, the real functions
\[
 f_\delta(x)=\delta\sin(x/\delta),\qquad\delta>0,
\]
converge uniformly to zero as $\delta\downarrow0$, but
$f_\delta'(0)=1$ for every $\delta$. The derivative of the limiting function
is zero. This is a real-family counterexample; no choice of a value for
$\sin\omega$ is needed to state it. Highly oscillatory or boundary-layer
families require additional estimates before reduction and differentiation
can be interchanged.

\subsection{Discarding the infinite part is not a canonical physical rule}

Consider the dimensionless Laurent expression
\[
 F(\eps)=\eps^{-1}+a,
 \qquad a\in\R.
\]
Its coefficient of $\eps^0$ is $a$. Under the equally legitimate formal
parameter change $\eps=\eta+c\eta^2$,
\begin{equation}
 F(\eta+c\eta^2)
   =\eta^{-1}+(a-c)+c^2\eta-\cdots.
 \label{eq:finitepart}
\end{equation}
The constant coefficient is now $a-c$. Thus ``keep the finite part'' for an
unlimited expression depends on a scale coordinate or subtraction scheme.
This does not contradict the uniqueness of standard part: $F$ was never in
$\mathcal O_F$.

Renormalization in a physical theory can include subtraction choices together
with covariance requirements and physical normalization conditions. Surreal
arithmetic does not make those requirements disappear. A coefficient-extraction
rule must be shown to respect the relevant coordinate and gauge symmetries,
and its observable consequences must not depend on arbitrary regulator
choices. Otherwise, a supposedly resolved curvature may be a scheme-dependent
number rather than a physical prediction.

\section{A regular-core example: why the dynamics, not the field, does the work}
\label{sec:Hayward}

Hayward's proposed nonsingular black-hole model provides a useful comparison
\cite{Hayward}. Consider its static metric function
\begin{equation}
 f_\ell(r)=1-\frac{2mr^2}{r^3+2m\ell^2},
 \qquad m>0,\quad\ell>0.
 \label{eq:hayward}
\end{equation}
For every fixed positive \emph{real} $\ell$, the central expansion is
\begin{equation}
 f_\ell(r)=1-\frac{r^2}{\ell^2}
                  +\frac{r^5}{2m\ell^4}+O(r^8).
 \label{eq:haywardexp}
\end{equation}
The core is de Sitter-like. Substitution into \eqref{eq:generalK} gives
\begin{equation}
 \Kretsch(0)=\frac{24}{\ell^4}.
 \label{eq:haywardK}
\end{equation}
For the metric convention used here,
\[
 G^t{}_t=\frac{r f'(r)+f(r)-1}{r^2}.
\]
With Einstein's equation $G^a{}_b=8\pi T^a{}_b$ and
$T^t{}_t=-\rho$, the core density is
\begin{equation}
 \rho(0)=\frac{3}{8\pi\ell^2}.
 \label{eq:haywardrho}
\end{equation}
These two central limits are checked by the companion program. Here a
``regular core'' refers to finite curvature and the corresponding classical
$C^2$-level metric regularity, not automatically to analytic or all-orders
smoothness. In fact, the $r^5$ term becomes $|x|^5$ along a Cartesian line
through the center, so this specific ansatz should not be advertised as a
$C^\infty$ Cartesian metric there. Unlike
Schwarzschild, the model is not vacuum in the core. Its effective stress tensor
and geometric modification are doing the regularizing work.

Now let $\ell=L\delta$ with $\delta$ a positive infinitesimal. The center has
\[
 L^4\Kretsch(0)=24\delta^{-4},
\]
which is unlimited. A family regular at every nonzero real regulator need not
have a bounded zero-regulator limit. Taking the regulator infinitesimal makes
that nonuniformity explicit; it does not supply a uniformly finite core.

The relevant crossover radius follows from balancing the denominator terms:
\begin{equation}
 r_{\mathrm c}=(2m\ell^2)^{1/3}.
 \label{eq:haywardcrossover}
\end{equation}
For $r^3\gg2m\ell^2$, the metric is approximately Schwarzschild. For
$r^3\ll2m\ell^2$, its leading central geometry is \eqref{eq:haywardexp}.
If $r/L=\eps^a$, $\ell/L=\eps^b$, and $m/L$ is ordinary and nonzero,
comparison of $3a$ with $2b$ identifies the outer and inner asymptotic sectors.
At equality one must retain both terms. This is another useful application of
ordered scale arithmetic and another warning against interchanging limits.

The regular center of this ansatz alone does not prove that a complete collapse
and evaporation model is globally stable, geodesically complete, or the
prediction of a consistent quantum theory. Those are additional physical and
geometric questions. General analyses of geodesically complete black-hole
models make such qualifications explicit \cite{CarballoRubio}. The lesson is
not that regular models are impossible, but that their success must be credited
to specified new geometry and dynamics rather than to allowing infinite
coefficients.

\section{Beyond Schwarzschild: weak singularities, Kerr, and oscillatory regimes}
\label{sec:beyond}

\subsection{Metric regularity and curvature regularity can separate}

The central Schwarzschild example is deliberately simple. Charged and rotating
interiors involve Cauchy horizons and instability mechanisms with different
regularity behavior. Dafermos and Luk prove a $C^0$ stability theorem for the
Kerr Cauchy horizon under their prescribed interior-data hypotheses, obtaining
continuous metric extensions across a nontrivial portion of that horizon
\cite{DafermosLuk}. This is not a theorem that generic black holes have a smooth,
unique, physically harmless interior continuation.

The result matters here because it demonstrates that even within real
mathematics, asking whether ``the singularity is removed'' without naming the
regularity class is incomplete. Finite or continuous metric components do not
imply bounded curvature, a $C^2$ extension, or uniqueness of the continued
Einstein evolution. Replacing those components by surreal numbers does not
settle any of these distinctions.

\subsection{Integrated tidal effects distinguish kinds of divergence}

Consider a positive tidal profile modeled near $s=0$ by
\begin{equation}
 T(s)=\frac{1}{s^2[\log(1/s)]^p},\qquad 0<s<s_0<1.
 \label{eq:logtidal}
\end{equation}
The logarithm uses a dimensionless time ratio; a reference time has been set
to one. As a simple diagnostic of twice-integrated tidal forcing, examine
\[
 \int_0^{s_0}sT(s)\,\dd s
 =\int_{\log(1/s_0)}^\infty u^{-p}\,\dd u.
\]
This is finite precisely when $p>1$. The weight $s$ arises by integrating a
positive forcing twice up to the endpoint and reversing the order of
integration. It is a useful model criterion, not a substitute for all
hypotheses in a geometric definition of singularity strength.

Thus two profiles that both have leading power $s^{-2}$ can have different
integrated effects because of logarithmic factors. A surreal/transseries
representation can preserve that distinction. The undifferentiated statement
``the tidal field is an infinite number'' cannot. Weak mass-inflation-type
singularities motivate examining such distinctions; representative charged
black-hole analyses explicitly separate diverging curvature from finite tidal
distortion \cite{Bhattacharjee}.

\subsection{Not every singular regime is a single tame transseries}

Another limitation is oscillation. Near a spacelike singularity, BKL-type
analyses describe changing anisotropic regimes, often represented by billiard
behavior. The scope and assumptions of such descriptions are nontrivial
\cite{DamourHenneauxNicolai}. One cannot infer that a generic dynamical interior
has one globally valid well-ordered power-log expansion merely because a
Schwarzschild geodesic has a Puiseux law.

Hardy fields and many asymptotic transseries constructions are designed for
eventually ordered, nonoscillatory behavior. Oscillatory phases, repeated
changes of dominant balance, or spatially localized features may require
additional variables, sectors, or a different function category. The real
surreal exponential does not by itself define an unrestricted periodic
function at every infinite argument. A successful application must prove that
the specific solution lies in its chosen class and that the representation
preserves the physical boundary or initial conditions.

\section{Surcomplex quantum mechanics: what survives, and what is missing}
\label{sec:quantum}

A proposal to replace both $\R$ and $\C$ should address quantum theory as well
as classical geometry. The answer is mixed: finite algebra survives very well;
the analytic and operational structures are not automatic.

\subsection{Finite-dimensional Hermitian algebra does survive}

Let $F$ be a real-closed ordered field containing $\R$, and let $K=F[\ii]$.
For $v,w\in K^n$, define
\[
 \langle v,w\rangle=\sum_{a=1}^n\overline{v_a}w_a.
\]
Then $\langle v,v\rangle>0$ for every $v\ne0$, because it is a finite sum of
squares of real and imaginary parts.

\begin{proposition}[Finite Hermitian spectral theorem]
\label{prop:spectral}
Every Hermitian $n\times n$ matrix over $K$ has an orthonormal basis of
eigenvectors and all its eigenvalues lie in $F$.
\end{proposition}
\begin{proof}
Since $K$ is algebraically closed, the characteristic polynomial has a root
$\lambda$ and there is a nonzero eigenvector $v$. Hermiticity gives
\[
 \lambda\langle v,v\rangle
 =\langle v,Av\rangle
 =\overline{\langle v,Av\rangle}
 =\overline\lambda\langle v,v\rangle.
\]
Positivity implies $\langle v,v\rangle\ne0$, so
$\lambda=\overline\lambda\in F$. The positive element
$\langle v,v\rangle$ has a square root, allowing normalization of $v$.
Its orthogonal complement is invariant under $A$, and induction gives an
orthonormal eigenbasis.
\end{proof}

This is a substantive algebraic starting point. It supports finite-state
amplitudes, projectors, observables, and unitary changes of basis over a
surcomplex workspace. Infinite eigenvalues can be represented; whether they
represent physical energies requires an interpretation and dynamics.

\subsection{Finite measurements have an ordinary shadow under a natural rule}

For a normalized state $\psi\in K^n$, finite-outcome Born probabilities are
\[
 p_a=|\psi_a|^2\in F,\qquad p_a\geq0,\qquad\sum_a p_a=1.
\]
Every component of $\psi$ is limited because its squared modulus is at most
one. Let $\psi_0=\st(\psi)$ componentwise. Then
\begin{equation}
 \sum_a|\psi_{0,a}|^2=1,
 \qquad \st(p_a)=|\psi_{0,a}|^2.
 \label{eq:bornshadow}
\end{equation}
Likewise, the entries of a finite unitary matrix are limited, and its standard
part is an ordinary complex unitary matrix. These statements follow directly
from the finite ring identities and positivity, not from an infinite-dimensional
spectral theorem.

Thus a model using ordinary finite apparatus and the standard-part readout
rule reproduces ordinary Born probabilities for these measurements. An
infinitesimal change need not be observable in any specified finite experiment
under that rule. This is not a theorem that every non-Archimedean model is
experimentally indistinguishable: dynamics or an amplification mechanism could
convert a small correction into a limited observable difference. Such a
mechanism and its finite resource scales must be provided. Invoking an
unlimited number of repetitions is not yet a finite experimental protocol.

\subsection{Probability requires a summation theory, not just positive scalars}

For infinitely many outcomes, the equation ``sum of probabilities equals one''
needs a definition. A countable family assigning the same positive
infinitesimal $\eps$ to each outcome is not strongly Hahn summable: infinitely
many summands contribute to the same exponent. Ordinary finite additivity
therefore cannot simply be promoted to a countably additive probability theory
by declaring each term infinitesimal.

Non-Archimedean probability theories exist and introduce explicit structures
for infinite addition and conditioning \cite{BenciHorstenWenmackers}. They are
examples of additional mathematics a physical proposal would have to choose,
not consequences of the fact that $F$ is an ordered field. Ordinary
measure-theoretic or nonstandard probability constructions can be incorporated
only after their relevant assumptions and interpretation are supplied.

\subsection{Time evolution and arbitrary infinite phases}

The expression
\[
 U(t)=\exp(-\ii Ht/\hbar)
\]
is not fully defined by the existence of $K=F[\ii]$. A real exponential on $F$
defines growth and decay, but it does not automatically choose values for
all oscillatory phases with unlimited arguments. A general complex exponential
requires a compatible function theory, periodic structure, and branch or
extension choices. Parametrizing the unit circle algebraically or using finite
angles does not by itself define a canonical value of $\sin\omega$.

There are tractable subclasses. One may keep the unperturbed real-time
propagator over $\C$ and solve a formal perturbation problem coefficientwise.
One may exponentiate a positive-order formal matrix, for which each coefficient
receives only finitely many relevant terms. One may work with an explicitly
constructed analytic extension on a restricted domain. What is not legitimate
is to infer a global unitary evolution for an unlimited Hamiltonian from the
finite Hermitian spectral theorem alone.

\subsection{Infinite-dimensional quantum theory and singular boundaries}

Quantum field theory and continuum quantum mechanics need more than finite
matrices: operator domains, infinite sums, spectral measures, distributions,
and suitable continuity or completeness. An algebraically closed scalar field
does not provide these structures, and a surreal-valued inner product is not
an ordinary Hilbert norm. A replacement theory must establish the relevant
functional analysis rather than silently import it.

There is also a distinct notion of quantum singularity. Horowitz and Marolf
study scalar quantum probes of static singular spacetimes and relate unique
quantum evolution to essential self-adjointness of the relevant spatial
operator \cite{HorowitzMarolf}. Their setting does not establish a universal
resolution of the spacelike Schwarzschild singularity. It does illustrate
what a meaningful question looks like: does the dynamics require extra
boundary conditions? Changing scalars alone does not select an operator domain
or remove an ambiguity in self-adjoint extension.

For semiclassical gravity, the problem is more demanding still. One must
specify the quantum state, renormalized stress tensor, and backreaction
prescription. None is selected by the availability of infinite surreal
energies or surcomplex amplitudes.

\section{The strongest positive connection: transseries, resurgence, and wave physics}
\label{sec:resurgence}

\subsection{Beyond-all-orders information is naturally ordered}

The hierarchy
\[
 0<e^{-1/\eps}\ll\eps^n\ll1
 \qquad(n\in\N\text{ fixed})
\]
expresses the fact that an exponentially small effect can be invisible to every
fixed algebraic order. Appropriate exponential transseries embed into the
surreal setting with compatible differential structure
\cite{BerarducciMantova,ADH}. This makes surreals a possible common language for
powers, logarithms, exponentials, and iterated scale hierarchies.

Such information is relevant to singular perturbation theory, tunneling,
semiclassical approximations, and wave matching. However, an asymptotic scale
is not itself a solution of a physical boundary-value problem. One must still
identify the actual function, determine free constants, and establish a valid
reconstruction or summation procedure.

\subsection{An explicit differential equation with hidden data}

Consider the dimensionless equation
\begin{equation}
 y'(x)+y(x)=\frac1x,
 \qquad x\to+\infty.
 \label{eq:factorialODE}
\end{equation}
A formal inverse-power solution is
\begin{equation}
 \widehat y(x)=\sum_{n\geq0}\frac{n!}{x^{n+1}}.
 \label{eq:factorialseries}
\end{equation}
To see this directly, substitute $\sum a_nx^{-n-1}$: $a_0=1$ and
$a_n=n a_{n-1}$, hence $a_n=n!$. For the truncated sum $y_N$,
\begin{equation}
 y_N'+y_N-\frac1x=-\frac{(N+1)!}{x^{N+2}}.
 \label{eq:factorialresidual}
\end{equation}
At $x=\omega$, \eqref{eq:factorialseries} is Hahn-summable. Under a compatible
strong derivation with $D\omega=1$, its differentiated terms cancel in the same
way. Nevertheless, the homogeneous equation has solutions $Ce^{-x}$. Thus
\begin{equation}
 y(x)=y_{\mathrm{part}}(x)+Ce^{-x}
 \label{eq:ambiguity}
\end{equation}
shares the same expansion in inverse powers for every ordinary constant $C$.
In the enlarged exponential field, $Ce^{-\omega}$ is a distinct,
beyond-all-orders scale.

The coefficient sequence in \eqref{eq:factorialseries} therefore does not
specify the complete solution. The value of $C$ must come from boundary or
initial data, or a prescribed summation and matching condition. A richer field
can \emph{retain} this information once supplied; it cannot infer it from data
that do not determine it. This is a useful miniature of the difficulty faced
when asymptotic black-hole solutions are mistaken for complete dynamics.

Costin and Ehrlich's integration theory is significant precisely because it
constructs controlled extensions for specified resurgent classes, rather than
asserting that every formal series automatically represents every desired
function \cite{CostinEhrlich}. Extending such methods to a particular reduced
gravitational system is a plausible project. Extending them to unrestricted
Einstein PDE dynamics would require substantially more analysis.

\subsection{A realistic black-hole benchmark}

Black-hole quasinormal-mode and scattering problems already use complex
analysis, asymptotic series, and boundary conditions. Hatsuda's work relates
black-hole quasinormal frequencies to high-order WKB expansions and shows, in
the cases studied, how Borel summation recovers the numerical frequencies
\cite{Hatsuda}. This provides a concrete benchmark for a surreal-compatible
implementation: preserve perturbative and nonperturbative sectors, enforce
the physical boundary conditions, and reproduce an established complex
frequency calculation.

Success would demonstrate calculational utility. It would not by itself
resolve an interior spacetime singularity, because the wave problem and its
boundary conditions are separate from a replacement of the central geometry.
The benchmark is attractive precisely because the claim is testable and
limited: compare coefficient generation, support certification, summation,
accuracy, and error control with established methods.

\subsection{Near-extremal and multiple-limit problems}

Another tractable algebraic example comes from the Reissner--Nordstr\"om
metric function $f(r)=1-2m/r+q^2/r^2$, whose horizon polynomial is
$r^2-2mr+q^2$. In geometrized units, let
$\delta=\sqrt{m^2-q^2}/m$ for $0<|q|<m$. Then
\[
 r_\pm=m(1\pm\delta),\qquad
 \kappa_+=\frac{r_+-r_-}{2r_+^2}
          =\frac{\delta}{m(1+\delta)^2}.
\]
A symbolic infinitesimal $\delta$ cleanly records the splitting of a double
root and the small surface-gravity scale. Multiple parameters can keep this
limit separate from a wave-frequency limit or a small perturbation amplitude.
These are algebraic identities and asymptotic bookkeeping, not claims that
extremality or a Cauchy horizon is physically harmless. The repository's
polynomial jets and finite-deformation work are much closer to this application
than to an already established non-Archimedean theory of spacetime.

\section{Comparison with alternative mathematical frameworks}
\label{sec:comparison}

A research program should explain why surreals are preferable to an existing
smaller tool for a particular task. Universality alone is not an efficiency or
prediction theorem.

\begin{center}
\small
\begin{tabular}{@{}P{0.22\textwidth}P{0.32\textwidth}P{0.35\textwidth}@{}}
\toprule
Framework & Natural advantage & What it does not supply by itself \\
\midrule
Formal power/Puiseux series & Finite coefficient algorithms; ramified scaling;
clear order-by-order equations & A convergent or uniquely selected real solution \\
Hahn fields & General ordered supports and multiple scales; exact algebraic
summation & Arbitrary smooth function theory or physical interpretation \\
Transseries and resurgence & Exponential sectors, asymptotic differential
structure, selected reconstruction methods & Universal applicability to nonlinear
spacetime PDEs \\
Hyperreal/nonstandard models & Infinitesimals with a separately constructed
transfer framework & New gravitational dynamics merely from changing scalars \\
Surreal/surcomplex framework & A broad common ambient field with algebraic,
exponential, and selected differential structures & Canonical topology, global
oscillatory calculus, probability, or singularity resolution \\
Weak or generalized geometric formulations & Explicitly altered regularity and
notions of solution & Uniqueness or physical admissibility without additional
conditions \\
\bottomrule
\end{tabular}
\end{center}

The comparison is structural, not a ranking of all possible theories. For an
individual expansion with rational powers, a Puiseux field may suffice. For
several ordered scales, a Hahn field may be the correct working object.
Surreals become attractive as an ambient framework connecting many such
structures or supporting a particularly useful normal form, not simply because
they are larger.

Likewise, representing a distributional source or multiplying singular fields
requires a theory of generalized functions and products. Substituting an
infinite scalar for a delta distribution does not retain its action on test
functions. A black-hole proposal involving weak stress tensors must address
that issue explicitly, irrespective of the chosen field.

\section{What the supplied repository contributes}
\label{sec:repo}

\subsection{Scope of the inspection}

The repository was inspected at commit
\begin{center}
\small\texttt{39f2be6667ade51bca2b45daa47e289d69c09764}.
\end{center}
The inspection covered the repository tree, root README, documentation catalog,
and surcomplex-analysis README. These are enough to assess the advertised
architecture and its stated limitations. This article does not claim a full
source audit, a new Lean build, or independent verification of all research
drafts \cite{RepoRoot,RepoDocs,RepoAnalysis}.

The catalog describes fifteen research packages. Its analytic reports
explicitly separate Hahn sums from topological limits, distinguish coefficient
rings and function classes, and warn that the reports are AI-assisted drafts
rather than fully refereed or machine-checked foundations. Those distinctions
are essential for any physics application, not incidental editorial details.

The root README describes implemented generic finite algebra, polynomial
operations and jets, complexification, and Hahn support and reduction results.
It explicitly states that these prerequisites do not yet construct the full
surreal field or establish its normal-form bridge to Hahn series. The source
documents have broader ambitions than the imported Lean statements. The
article's physical conclusions therefore do not rest on a claim that the
repository already formalizes surcomplex analysis, Lorentzian geometry, or
Einstein evolution in full \cite{RepoRoot}.

\subsection{The most relevant existing ingredients}

The finite algebra layer is relevant to polynomial horizon equations,
degenerating eigenvalue problems, tensor contractions, and finite jets. The
Hahn layer's explicit support conditions are relevant to perturbative
inversion, composition, and regrouping. Coefficient-zero reduction is relevant
to a formal correspondence map such as Proposition~\ref{prop:reduction}.

The analytic documents' common-domain warnings are especially important for
physics. A sequence of coefficient functions that are each analytic somewhere
near a point need not share a neighborhood on which an infinite coefficient
family defines the intended object. Uniform domain and support assumptions
must accompany any claim about a perturbative metric, wave solution, or
integration procedure \cite{RepoDocs,RepoAnalysis}.

The trigonometric and complex-algebraic material should not be read as a
finished unrestricted phase calculus for quantum evolution. Similarly, a
polynomial factorization theorem is not an existence theorem for a nonlinear
hyperbolic PDE. Recognizing these boundaries allows the repository to be used
productively without overinterpreting it.

\subsection{A focused extension plan}

A first physics-oriented layer should be generic over a controlled ordered
coefficient field or formal differential algebra, rather than require all of
$\No$ at the outset. The following are proposed additions, not descriptions of
existing modules.


\paragraph{Formal metric algebra.}
Represent symmetric matrices with an invertible leading metric, prove the
formal inverse identity, and implement connection and curvature constructions
with coefficientwise derivatives.
\paragraph{Regular reduction and gauge covariance.}
Formalize Proposition~\ref{prop:reduction}, the finite Bianchi identities, and
compatibility with ordinary coordinate transformations. Record the exact
coefficient-ring assumptions.
\paragraph{Certified benchmark calculations.}
Reproduce the Eddington--Finkelstein determinant, Schwarzschild curvature,
proper-time and Jacobi exponents, angular crossover balances, and Hayward core
limits. Distinguish symbolic identity checking from analytic error estimates.
\paragraph{Asymptotic realization interfaces.}
For a selected ODE or wave problem, attach a real-domain solution theorem,
remainder bounds, and a summation or matching prescription to the formal
series. Do not expose an unrestricted evaluation operation lacking hypotheses.
\paragraph{Operational correspondence.}
Specify which quantities are limited, when standard part commutes with an
operation, and which outputs are intended as physical observables.


This route produces reusable mathematics even if surreals never become
fundamental physical scalars. It also creates explicit interfaces where a
future genuinely non-Archimedean physical proposal would have to add new
assumptions.

\section{Published surreal-physics proposals and the status of the idea}
\label{sec:literaturestatus}

There are published and preprint discussions connecting surreals to physics.
Nieto discusses links involving dualities, matroids, quantum structures, and
gravity, and suggests that available surreal infinities may help with
black-hole or cosmological divergences \cite{Nieto2016,Nieto2018}. Such work
establishes that the suggestion has precedents; it does not justify dismissing
it as an unasked or meaningless question.

The specific inference from ``infinite values belong to the number system''
to ``the singularity problem has been solved'' needs more argument.
Equations~\eqref{eq:infK}, \eqref{eq:radial}, and \eqref{eq:polyobstruction}
exhibit what is missing: unlimited invariant curvature remains unlimited, a
finite proper-time endpoint is still present, and the exact central identity
still has no field-valued solution. A complete alternative might replace the
metric, state space, observables, or dynamical law, but those replacements
must be constructed.

The literature reviewed here does not establish an empirically tested theory
in which the scalar replacement alone resolves a black-hole singularity.
This is a bounded assessment of the cited work, not a proof that no future
non-Archimedean theory can do so. Conversely, the established mathematical
results on surreal differential fields and integration are genuine and should
not be conflated with speculative gravitational interpretations
\cite{BerarducciMantova,ADH,CostinEhrlich}.

\section{A research program with clear success criteria}
\label{sec:research}

The central research challenge is to make a proposal precise enough that it
can fail a test. A larger collection of named numbers is not yet such a
proposal.

\subsection{A minimal model specification}

A candidate theory should specify its scalar workspace and foundational size;
its spaces of states, events, and fields; its topology and derivatives; its
rules for infinite sums and integration; its field equations and admissible
initial data; its causal or quantum evolution law; its gauge and unit
conventions; and its operational readout. Recovery of ordinary physics must
be a theorem or a controlled approximation in the chosen model, not an appeal
to the inclusion $\R\subset\No$.

For a singularity claim, the specification must additionally state what
``resolved'' means. Examples include extension of every relevant causal
geodesic, bounded detector responses, existence and uniqueness of weak
evolution, or a specified quantum propagation criterion. These conditions
are not equivalent. A model may intentionally replace classical geodesics,
but then it must replace the question of their termination with an operationally
meaningful one.

\subsection{Projects that are realistic and scientifically useful}

\paragraph{Project A: certified multiscale infall.}
Develop the crossover problem \eqref{eq:crossoverexact} in a controlled
Puiseux/Hahn setting, prove uniform remainder estimates on matching domains,
and compare with high-precision real integration. A success is a theorem
connecting the symbolic sectors to actual geodesics, not merely a list of
formal dominant terms.

\paragraph{Project B: a regular-core singular perturbation problem.}
Use \eqref{eq:hayward} or another explicitly specified model to track the
outer Schwarzschild and inner-core expansions, their overlap, and invariant
curvature bounds as functions of the regulator. A success distinguishes
regularity at fixed regulator from uniform control as the regulator vanishes,
and identifies the required stress tensor.

\paragraph{Project C: a wave/resurgence benchmark.}
Select a single black-hole wave equation with known boundary conditions and
well-tested frequencies. Construct its transseries representation, prove the
needed support and summation statements, and recover the physical solution.
A success is reproducible accuracy or a new theorem about the asymptotic
structure, not a numerical coincidence from an unspecified infinite sum.

\paragraph{Project D: a genuine non-Archimedean modification.}
Only after the preceding interfaces are understood, choose an explicit new
dynamical or operational postulate. Prove a well-posedness result in a reduced
model, show what replaces the classical endpoint, and determine whether the
ordinary observable prediction differs from the corresponding real model.
A success would be a substantive physical proposal; no such postulate is
forced by surreal arithmetic alone.

\subsection{What would constitute a decisive advance?}

On the mathematical side, an extension theorem for a relevant gravitational
solution class, compatible with gauge transformations and controlled remainder
bounds, would be valuable. So would a rigorous matching theorem across a
singular perturbation layer or a spectral result for a precisely defined
surcomplex wave operator.

On the physical side, a decisive advance would require a replacement evolution
that remains predictive where the classical model ceases to be, together with
an interpretation of its observables and a demonstrable ordinary limit.
An observable consequence would have to come from this additional dynamics or
interpretation, rather than from the verbal reclassification of a divergent
quantity as an infinite field element.

\section{Conclusion}
\label{sec:conclusion}

Surreal and surcomplex numbers can represent the hierarchy of quantities near
a singular regime with exceptional algebraic precision. The explicit examples
in this article preserve powers of proper time, competing angular-momentum
scales, logarithmic corrections, and exponentially small sectors. Established
surreal differential and integration theories make selected applications more
than speculative notation.

They do not make the original Schwarzschild singularity regular simply by
admitting its nearby infinite values. Its invariant curvature, clock behavior,
and tidal equations survive the scalar extension. The exact center still
requires a new geometric or physical construction. Nor can one replace $\R$
and $\C$ wholesale while leaving the topology, differential equations,
probability, and quantum functional analysis unspecified.

The most productive interpretation is therefore \emph{surreal-compatible
asymptotic physics}: retain ordinary spacetime and operational observables,
use a carefully chosen formal or Hahn coefficient structure, connect it to
actual solutions by proved realization and matching results, and formalize
the algebraic interfaces. A more radical theory remains possible, but its
singularity-resolving content must lie in explicit new dynamics or a rigorously
defined replacement of the classical state space. The larger field can help
express that theory; it cannot supply it on its own.

\appendix
\section{Curvature calculation and dimensional conventions}
\label{app:curvature}

For clarity, the curvature convention used by the companion checker is
\[
 R^a{}_{bcd}=\partial_c\Gamma^a_{db}-\partial_d\Gamma^a_{cb}
       +\Gamma^a_{ce}\Gamma^e_{db}-\Gamma^a_{de}\Gamma^e_{cb},
 \qquad R_{bd}=R^a{}_{bad}.
\]
With $g=\diag(-f,f^{-1},r^2,r^2\sin^2\theta)$, the independent curvature
magnitudes in an orthonormal description are built from $f''/2$, $f'/(2r)$,
and $(1-f)/r^2$. Their multiplicities yield
\[
 R_{abcd}R^{abcd}=(f'')^2+\frac{4(f')^2}{r^2}
                        +\frac{4(1-f)^2}{r^4}.
\]
The coordinate calculation in the script checks the same result without
relying on a static orthonormal frame across the horizon. It is carried out
where the original chart is nonsingular and gives an invariant rational
expression that can then be evaluated in a horizon-regular chart.

All curvature powers are stated with $[\Kretsch]=\text{length}^{-4}$. The
proper-time variable in equations such as \eqref{eq:radial} is a length;
if $\Delta\tau_{\rm sec}$ is measured in seconds, replace $s$ by
$c\Delta\tau_{\rm sec}$. For example,
\[
 \Kretsch=\frac{64}{27c^4(\Delta\tau_{\rm sec})^4}
\]
for the specified $E=1$ radial infall. This convention prevents dimensional
constants from being mistaken for new asymptotic exponents.

The scalar $\Kretsch$ is a particularly convenient diagnostic for the
Schwarzschild and regular-core examples, where it is positive. It is not a
universal positive norm of curvature in Lorentzian geometry. More general
spacetimes can require frame components or other invariants to detect a
singular regime. No conclusion in this article equates bounded Kretschmann
scalar with general spacetime regularity.

\section{Verification and reproducibility}
\label{app:verification}

The archive includes \texttt{verify\_examples.py} and its recorded output.
The program uses SymPy to compute the Levi-Civita connection and full Riemann
tensor of the generic spherical metric, contract its invariant, and verify
selected identities used in the article. It then checks the Schwarzschild
Ricci tensor, the regular-horizon determinant, proper-time curvature, Jacobi
indicial solutions, the nonradial dominant balance, the exact crossover
integrand factor, the regular-core curvature and density limits, the
finite-part reparameterization, and factorial-series ODE residuals at eight
finite truncation orders.

The recorded run passed all \textbf{37 checks} under Python 3.13.5 and
SymPy 1.14.0. The tests are finite symbolic calculations. They do not verify
Propositions~\ref{prop:discrete}, \ref{prop:reduction}, or
\ref{prop:spectral} in a proof assistant; these have ordinary mathematical
proofs in the text. Nor do the tests establish convergence of a transseries,
validate a physical model, or execute the repository's Lean project.

The LaTeX source uses an internal bibliography and no external graphics.
It can be built using two or more runs of \texttt{pdflatex}, or with
\texttt{latexmk -pdf surreal\_physics.tex}. The archive README records the
build and the scope of verification. Bibliographic references distinguish
peer-reviewed mathematical results, physics papers, speculative connections,
and repository documentation.

\section{Compact claim ledger}
\label{app:ledger}

\begin{longtable}{@{}P{0.43\textwidth}P{0.46\textwidth}@{}}
\toprule
Claim & Status and dependency \\
\midrule
\endfirsthead
\toprule
Claim & Status and dependency \\
\midrule
\endhead
Surreal/surcomplex algebra supports ordered and algebraically closed
extensions. & Established mathematical background; not a complete physical
model \cite{Conway,Gonshor}. \\
Full order-topological real histories are constant. & Proved here for ordinary
connected real domains by set-sized discreteness; other topologies are not
excluded. \\
Surreal calculus has substantial rigorous subclasses. & Established results
on differential fields and resurgent integration; no general Einstein PDE
claim \cite{BerarducciMantova,ADH,CostinEhrlich}. \\
Schwarzschild $r=L\eps$ has unlimited curvature. & Exact invariant calculation;
verified symbolically and explained by Proposition~\ref{prop:pole}. \\
Radial $E=1$ infall has $\Kretsch=64/(27s^4)$. & Exact calculation for the named
geodesic class; not the universal nonradial law. \\
Small angular momentum produces a crossover. & Exact integral identity plus
stated bounded-domain limit and asymptotic balances; no new dynamics. \\
Regular formal Einstein solutions reduce to real ones. & Proved in the
specified differential algebra with invertible smooth leading metric. \\
A finite real core parameter can regularize the model's center. & Explicit
Hayward metric calculation; global physical viability is a separate issue. \\
Finite surcomplex Hermitian quantum algebra works. & Proved for finite matrices
over the complexification of a real-closed field. \\
The scalar replacement alone resolves physical black-hole singularities. & Not
established by the reviewed sources; the unchanged Schwarzschild identities
show why assigning infinite values is insufficient. \\
A future non-Archimedean theory might yield new physics. & Open possibility,
conditional on explicit dynamics, analytic structure, and observational
interpretation; not a result claimed here. \\
\bottomrule
\end{longtable}

\clearpage
\phantomsection
\addcontentsline{toc}{section}{References}
\begingroup
\raggedright
\begin{thebibliography}{99}
\setlength{\itemsep}{0.45em}

\bibitem{Conway}
J. H. Conway, \emph{On Numbers and Games}, second edition,
A K Peters, 2001; first edition, Academic Press, 1976.

\bibitem{Gonshor}
H. Gonshor, \emph{An Introduction to the Theory of Surreal Numbers},
London Mathematical Society Lecture Note Series 110,
Cambridge University Press, 1986.

\bibitem{vanDenDriesEhrlich}
L. van den Dries and P. Ehrlich,
``Fields of surreal numbers and exponentiation,''
\emph{Fundamenta Mathematicae} \textbf{167} (2001), 173--188.
\href{https://doi.org/10.4064/fm167-2-3}{doi:10.4064/fm167-2-3}.

\bibitem{BerarducciMantova}
A. Berarducci and V. Mantova,
``Surreal numbers, derivations and transseries,''
\emph{Journal of the European Mathematical Society} \textbf{20} (2018), 339--390.
\href{https://doi.org/10.4171/JEMS/769}{doi:10.4171/JEMS/769};
\href{https://arxiv.org/abs/1503.00315}{arXiv:1503.00315}.

\bibitem{ADH}
M. Aschenbrenner, L. van den Dries, and J. van der Hoeven,
``The surreal numbers as a universal $H$-field,''
\href{https://arxiv.org/abs/1512.02267}{arXiv:1512.02267},
version 3, 2016. The cited version records acceptance for publication in the
\emph{Journal of the European Mathematical Society}.

\bibitem{CostinEhrlich}
O. Costin and P. Ehrlich,
``Integration on the surreals,''
\emph{Advances in Mathematics} \textbf{452} (2024), 109823.
\href{https://arxiv.org/abs/2208.14331}{arXiv:2208.14331};
\href{https://arxiv.org/html/2208.14331v5}{full text, version 5}.

\bibitem{Penrose}
R. Penrose, ``Gravitational collapse and space-time singularities,''
\emph{Physical Review Letters} \textbf{14} (1965), 57--59.
\href{https://doi.org/10.1103/PhysRevLett.14.57}{doi:10.1103/PhysRevLett.14.57}.

\bibitem{Witten}
E. Witten, ``Light rays, singularities, and all that,''
\emph{Reviews of Modern Physics} \textbf{92} (2020), 045004.
\href{https://doi.org/10.1103/RevModPhys.92.045004}{doi:10.1103/RevModPhys.92.045004};
\href{https://arxiv.org/abs/1901.03928}{arXiv:1901.03928}.

\bibitem{Landsman}
K. Landsman, ``Penrose's 1965 singularity theorem: From geodesic incompleteness
to cosmic censorship,''
\href{https://arxiv.org/abs/2205.01680}{arXiv:2205.01680}, 2022.

\bibitem{DafermosLuk}
M. Dafermos and J. Luk,
``The interior of dynamical vacuum black holes I: The $C^0$-stability of the
Kerr Cauchy horizon,''
\emph{Annals of Mathematics} \textbf{202} (2025), 309--630.
\href{https://doi.org/10.4007/annals.2025.202.2.1}{doi:10.4007/annals.2025.202.2.1};
\href{https://arxiv.org/abs/1710.01722}{arXiv:1710.01722}.

\bibitem{Donoghue}
J. F. Donoghue,
``General relativity as an effective field theory: The leading quantum corrections,''
\emph{Physical Review D} \textbf{50} (1994), 3874--3888.
\href{https://doi.org/10.1103/PhysRevD.50.3874}{doi:10.1103/PhysRevD.50.3874};
\href{https://arxiv.org/abs/gr-qc/9405057}{arXiv:gr-qc/9405057}.

\bibitem{Hayward}
S. A. Hayward, ``Formation and evaporation of non-singular black holes,''
\emph{Physical Review Letters} \textbf{96} (2006), 031103.
\href{https://doi.org/10.1103/PhysRevLett.96.031103}{doi:10.1103/PhysRevLett.96.031103};
\href{https://arxiv.org/abs/gr-qc/0506126}{arXiv:gr-qc/0506126}.

\bibitem{CarballoRubio}
R. Carballo-Rubio, F. Di Filippo, S. Liberati, and M. Visser,
``Geodesically complete black holes,''
\emph{Physical Review D} \textbf{101} (2020), 084047.
\href{https://doi.org/10.1103/PhysRevD.101.084047}{doi:10.1103/PhysRevD.101.084047};
\href{https://arxiv.org/abs/1911.11200}{arXiv:1911.11200}.

\bibitem{Bhattacharjee}
S. Bhattacharjee, S. Sarkar, and A. Virmani,
``Internal structure of charged AdS black holes,''
\href{https://arxiv.org/abs/1604.03730}{arXiv:1604.03730}, 2016.

\bibitem{DamourHenneauxNicolai}
T. Damour, M. Henneaux, and H. Nicolai,
``Cosmological billiards,''
\emph{Classical and Quantum Gravity} \textbf{20} (2003), R145--R200.
\href{https://arxiv.org/abs/hep-th/0212256}{arXiv:hep-th/0212256}.

\bibitem{BenciHorstenWenmackers}
V. Benci, L. Horsten, and S. Wenmackers,
``Non-Archimedean probability,''
\emph{Milan Journal of Mathematics} \textbf{81} (2013), 121--151.
\href{https://arxiv.org/abs/1106.1524}{arXiv:1106.1524}.

\bibitem{HorowitzMarolf}
G. T. Horowitz and D. Marolf,
``Quantum probes of spacetime singularities,''
\emph{Physical Review D} \textbf{52} (1995), 5670--5675.
\href{https://doi.org/10.1103/PhysRevD.52.5670}{doi:10.1103/PhysRevD.52.5670};
\href{https://arxiv.org/abs/gr-qc/9504028}{arXiv:gr-qc/9504028}.

\bibitem{Hatsuda}
Y. Hatsuda, ``Quasinormal modes of black holes and Borel summation,''
\emph{Physical Review D} \textbf{101} (2020), 024008.
\href{https://doi.org/10.1103/PhysRevD.101.024008}{doi:10.1103/PhysRevD.101.024008};
\href{https://arxiv.org/abs/1906.07232}{arXiv:1906.07232}.

\bibitem{Nieto2016}
J. A. Nieto, ``Some mathematical and physical remarks on surreal numbers,''
2016.
\href{https://doi.org/10.4236/jmp.2016.715188}{doi:10.4236/jmp.2016.715188};
\href{https://arxiv.org/abs/1611.09699}{arXiv:1611.09699}.

\bibitem{Nieto2018}
J. A. Nieto, ``Duality, matroids, qubits, twistors and surreal numbers,''
\href{https://arxiv.org/abs/1810.04521}{arXiv:1810.04521}, 2018.

\bibitem{RepoRoot}
\texttt{VladimirReshetnikov/Surreal}, root \texttt{README.md}.
\href{https://github.com/VladimirReshetnikov/Surreal/blob/39f2be6667ade51bca2b45daa47e289d69c09764/README.md}{Pinned repository overview},
commit \texttt{39f2be6667ade51bca2b45daa47e289d69c09764},
accessed 21 September 2026. Scope and limitations of implemented Lean prerequisites.

\bibitem{RepoDocs}
\texttt{VladimirReshetnikov/Surreal}, \texttt{docs/README.md}.
\href{https://github.com/VladimirReshetnikov/Surreal/blob/39f2be6667ade51bca2b45daa47e289d69c09764/docs/README.md}{Pinned documentation catalog},
same commit and access date. Includes the explicit status and coefficient-category
warnings for the research reports.

\bibitem{RepoAnalysis}
\texttt{VladimirReshetnikov/Surreal},
\texttt{docs/surcomplex/analysis/README.md}.
\href{https://github.com/VladimirReshetnikov/Surreal/blob/39f2be6667ade51bca2b45daa47e289d69c09764/docs/surcomplex/analysis/README.md}{Pinned analysis overview},
same commit and access date. Distinguishes summability, function classes, and
residue notions; not treated here as a fully verified physics foundation.

\end{thebibliography}
\endgroup
\end{document}
